\chapter{调度器深入：时间、等待与多核}

\section{原理}

上一章讲了基本调度算法，本章把调度器当成内核子系统看。调度器不是一个简单的 \code{pick\_next()} 函数，它维护可运行队列、睡眠队列、时间片账本、优先级、CPU 亲和性、负载均衡、抢占标志和统计数据。它还必须和锁、定时器、中断、系统调用返回路径配合。

调度器的核心循环是：当前任务运行；某个事件改变任务状态；如果当前任务不能继续运行或应该让出 CPU，内核选择下一个 runnable 任务；保存当前上下文，恢复下一个上下文。事件可能来自时间片耗尽、任务阻塞、任务退出、任务被唤醒、高优先级任务到达或 CPU 负载均衡。

调度器维护两个不同问题。第一是选择：在 runnable 集合里选谁。第二是状态转换：任务如何进入 runnable、running、sleeping、stopped、zombie。很多 bug 不在选择算法，而在状态转换，例如任务已经 sleeping 却仍留在 runqueue，或者唤醒丢失导致任务永远睡眠。

\section{实践}

实践报告要把调度器拆成四张表：

\begin{longtable}{p{0.2\linewidth}p{0.68\linewidth}}
\toprule
表 & 内容 \\
\midrule
任务表 & pid/tid、状态、优先级、时间片、等待原因 \\
runqueue & 每个 CPU 的 runnable 任务和顺序 \\
sleep queue & 等待对象、超时时间、唤醒条件 \\
统计表 & on-CPU 时间、runnable wait、context switch、迁移次数 \\
\bottomrule
\end{longtable}

当报告说“调度不公平”时，必须给出 runnable wait 分布；当报告说“线程没运行”时，必须说明它是 runnable 但没拿到 CPU，还是 sleeping 等事件。

\section{算法与实现分析}

\subsection{睡眠和唤醒}

阻塞系统调用和条件变量最终都要把任务放入某个等待队列。正确顺序是：持有保护等待条件的锁，检查条件，不满足则把当前任务加入等待队列并把状态改为 sleeping，然后原子地释放锁并调度。唤醒方改变条件后，从等待队列取任务，改为 runnable，放入 runqueue。

\begin{lstlisting}
sleep(wait_queue, lock):
  current.state = SLEEPING
  wait_queue.push(current)
  unlock(lock)
  schedule()
  lock(lock)

wake_one(wait_queue):
  task = wait_queue.pop()
  if task != null:
    task.state = RUNNABLE
    enqueue(runqueue_for(task), task)
\end{lstlisting}

这里最危险的是 lost wakeup。如果等待者检查条件后、真正入队前，唤醒者改变条件并发出唤醒，那么等待者可能永远睡眠。解决方法是用同一把锁保护“检查条件”和“加入等待队列”。

\subsection{抢占点}

内核不能在任意机器指令之间随便切换。抢占点通常出现在中断返回、系统调用返回、显式阻塞、时间片耗尽或内核允许抢占的位置。如果当前代码持有自旋锁或处于中断上下文，就不能睡眠调度。

\begin{lstlisting}
maybe_preempt():
  if need_reschedule
     and preempt_count == 0
     and !in_interrupt():
       schedule()
\end{lstlisting}

\code{preempt\_count} 是实现边界的典型例子。它不是业务概念，而是内核用来防止在不可抢占区域切走任务的计数。进入自旋锁或中断上下文时增加，退出时减少。若计数错误，可能在持锁时调度，导致死锁。

\subsection{每 CPU runqueue}

多核系统若使用一个全局 runqueue，所有 CPU 都要争同一把锁。每 CPU runqueue 降低锁竞争：每个 CPU 主要从自己的队列取任务；周期性负载均衡把任务从繁忙 CPU 迁到空闲 CPU。

\begin{lstlisting}
schedule_on_cpu(cpu):
  rq = cpu.runqueue
  lock(rq.lock)
  next = pick_next(rq)
  unlock(rq.lock)
  switch_to(next)

load_balance():
  busiest = find_busiest_cpu()
  idle = find_idle_cpu()
  migrate_some_tasks(busiest.rq, idle.rq)
\end{lstlisting}

多核调度的取舍是局部性和均衡。任务留在原 CPU，cache/TLB 更热；迁移到空闲 CPU，等待时间更低。调度器不能只看队列长度，还要看任务亲和性、NUMA 节点、迁移成本和负载类型。

\subsection{实时任务}

实时调度关注 deadline，而不只是平均公平。Rate Monotonic 适合周期任务，周期越短优先级越高；Earliest Deadline First 总是选择绝对 deadline 最近的任务。单处理器上，在理想条件下 EDF 可以达到更高利用率，但实现要维护按 deadline 排序的队列。

\begin{lstlisting}
pick_edf(runqueue):
  return min_by(runqueue, task.absolute_deadline)
\end{lstlisting}

实时系统必须做可调度性分析。若任务集合的最坏执行时间总和超过 CPU 能力，任何调度算法都无法保证 deadline。操作系统只能执行策略，不能创造时间。

\subsection{唤醒抢占}

一个高优先级或交互任务被 I/O 唤醒时，如果只是放到 runqueue 尾部，可能仍要等当前 CPU 密集任务用完整个时间片。唤醒抢占的思想是：唤醒路径比较新任务和当前任务，若新任务更应该运行，就设置 \code{need\_reschedule}，让返回路径尽快切换。

\begin{lstlisting}
try_wake_task(task):
  rq = select_runqueue(task)
  lock(rq.lock)
  task.state = RUNNABLE
  enqueue(rq, task)
  if should_preempt(task, rq.current):
    rq.current.need_reschedule = true
    send_reschedule_ipi_if_remote(rq.cpu)
  unlock(rq.lock)
\end{lstlisting}

这里有两个风险。第一，过度抢占会增加上下文切换，降低吞吐。第二，跨 CPU 唤醒需要 IPI，若频率过高，会把多核机器变成互相打断的系统。调度器必须在响应延迟和扰动成本之间折中。

\subsection{负载均衡的粒度}

负载均衡不是简单平均任务数量。一个 CPU 上有 1 个 CPU 密集任务，另一个 CPU 上有 20 个大部分睡眠的任务，按数量看不平衡，按实际负载可能合理。更好的指标是 runnable load、最近 CPU 使用、任务权重和亲和性。

\begin{lstlisting}
load_balance_domain(domain):
  busiest = cpu with highest runnable_load
  idle = cpu with lowest runnable_load
  if busiest.load <= idle.load + imbalance_threshold:
    return
  task = choose_migratable_task(busiest, idle)
  if task != null:
    migrate(task, busiest, idle)
\end{lstlisting}

迁移必须持有源和目标 runqueue 锁，并规定锁顺序，例如按 CPU id 从小到大加锁。否则两个 CPU 同时互相迁移任务，会在 runqueue 锁上形成死锁。

\subsection{调度器不变量}

调度器测试要检查不变量，而不是只看任务有没有输出。

\begin{rubricbox}
\begin{itemize}
\tightlist
\item 每个 RUNNABLE 线程恰好在一个 runqueue 中。
\item RUNNING 线程只能是某个 CPU 的 current。
\item SLEEPING 线程必须在一个等待队列中，或带有可解释超时。
\item ZOMBIE 进程不能再获得 CPU，但退出状态必须可被父进程读取。
\item 同一线程不能同时出现在迁移源队列和目标队列。
\item 持有自旋锁、关中断或 \code{preempt\_count > 0} 时不能调度到会睡眠的路径。
\end{itemize}
\end{rubricbox}

\section{测试}

\begin{itemize}
\tightlist
\item lost wakeup 测试：唤醒和睡眠并发发生时，等待者不会永久睡眠。
\item 持有自旋锁时不能进入会睡眠的路径。
\item 每 CPU runqueue 中，同一任务不能同时出现在两个队列。
\item 任务迁移后，旧 CPU 不再能调度该任务。
\item 时间片耗尽、阻塞、退出和唤醒都能触发正确重调度。
\item 实时任务 deadline miss 必须被记录，而不是静默忽略。
\end{itemize}
